% ===================================================================
%  TABELO N° 4 — Matur-evo ed oldeso.
%  Feuillets 26 a 31 du fac-simile = folios 24 a 29.
%  Correspondance : folio imprime = numero de feuillet - 2.
%
%  Ca tabelo esas partigita per l'autoro en du « SCENI » (Unesma
%  ceno, Duesma ceno) : singla sceno riincepas sua propra numerado
%  di alinei ye 1. La kloci %%K sequas la konsigno — la numero
%  imprimita da l'autoro — e do YE 1 REKOMENCAS por la Duesma ceno :
%  t04-01-1 aparas do DU FOYE en ca file, unfoye por singla sceno.
%  To ne esas eroro di transskripto : olu reflektas fidele la
%  numerado dil fac-simile.
% ===================================================================

% --- Feuillet 26 (folio 24) -----------------------------------------
\begin{VUpage}[26]{24}
%%K t04-apar-1 apar
\VUsaut{4.0mm}
\VUtitre{15.0pt}{60}{TABELO N\textsuperscript{o} 4}
\VUsaut{1.6mm}
\VUtitre{11.4pt}{0}{Matur - evo ed oldeso.}
\VUsaut{3.2mm}

%%K t04-tit-1 sub
\VUcentre{11.4pt}{0}{\textit{Unesma ceno.}}
\VUsaut{1.6mm}
\VUtitre{11.4pt}{0}{La Mariaj-festino.}
\VUsaut{2.6mm}

%%K t04-tit-2 sub
\VUpk{10.2pt}{40}{L'Autuno.}
\VUsaut{3.0mm}

%%K t04-c1-01-1 p
\VUblancAlinea
1. --- La yuni plugrandeskis. Un de li, Ioan\cc
nes, su mariajas. Ni asistas la \VUgras{mariaj-festino}\nl
qua asemblas, cirkum la sama tablo, la familii\nl
di la gespozi.

%%K t04-c1-02-1 p
\VUblancAlinea
2. --- Ni esas en \VUgras{autuno,} en la \VUgras{septembro}\cc
\VUgras{monato.} La kolda vento di \VUgras{oktobro} e di \VUgras{no}\cc
\VUgras{vembro} ne ja spoliis la ruro de lua verda fo\cc
liaro. La vesperal aero esas tepida e parfu\cc
moza; pro to la dineo eventas sub la \VUgras{verando} \textsuperscript{(8)}\nl
an la gardeno.

%%K t04-c1-03-1 p
\VUblancAlinea
3. --- Fore, sur la \VUgras{kolino} \textsuperscript{(3)}, situesas la hemo\nl
di la gepatri di Ioannes, eleganta \VUgras{kastelo} \textsuperscript{(1)}\nl
celata en la verdajo; bela viteyi, qui produk\cc
tas bonega vino, kovras la pento di la koli\cc
neto. La vitkultivisto kultivas e sorgas li.

%%K t04-c1-04-1 p
\VUblancAlinea
4. --- Super la viteyi, pasas flugo de \VUgras{per}\cc
\VUgras{driki} \textsuperscript{(2)} pavorigita dal proximesko dil \VUgras{chas}\cc
\VUgras{gardisto} \textsuperscript{(5)}. Ilca, kun \VUgras{pafilo} sub brakio, e ban\cc
\VUgras{doliera vildo-sako} exploras la \VUgras{busharo} \textsuperscript{(4)} kun\nl
sua \VUgras{hundo} \textsuperscript{(6)}. Il surveyas la domeno ed impe\cc
das la furte-chaseri destruktar la vildo.

%%K t04-c1-05-1 p
\VUblancAlinea
5. --- Exter la \VUgras{verando,} klimas \VUgras{vit-espaliero}\nl
de qua pendas densa \VUgras{vito-grapi} \textsuperscript{(7)}.
\end{VUpage}

% --- Feuillet 27 (folio 25) -----------------------------------------
\begin{VUpage}[27]{25}
%%K t04-c1-06-1 p
\VUblancAlinea
6. --- La bela flori autunala, qui garnisas la\nl
\VUgras{tablo} \textsuperscript{(16)}, esas \VUgras{krizantemi} \textsuperscript{(9)}; la \VUgras{patrulo} \textsuperscript{(10)} di\nl
la mariajitino metis un de oli en la butontruo\nl
di sua frako.

%%K t04-c1-07-1 p
\VUblancAlinea
7. --- La repasto fineskas. La \VUgras{servisto} \textsuperscript{(11)} ko\cc
mencas adportar la deser-dishi e balde depre\cc
nos la diversa parti dil manjilaro, \VUgras{kulieri} \textsuperscript{(14)},\nl
\VUgras{forketi} \textsuperscript{(12)}, \VUgras{kulteli} \textsuperscript{(13)}, \VUgras{pladegi} \textsuperscript{(15)} quin on ne\nl
plus bezonas.

%%K t04-tit-3 sub
\VUpk{10.2pt}{40}{La Parentaro.}
\VUsaut{3.0mm}

%%K t04-c1-08-1 p
\VUblancAlinea
8. --- Omni aspektas joyoza, e la rideto kun\nl
qua la \VUgras{matro} \textsuperscript{(17)} dil mariajitulo regardas sua\nl
filii pruvas lua feliceso.

%%K t04-c1-09-1 p
\VUblancAlinea
9. --- Ta mariajo unionas du richa familii di\nl
la lando. Ioannes, la \VUgras{spozulo} \textsuperscript{(18)}, esas \VUgras{filio} di\nl
granda proprietero ed il divenas la \VUgras{bofilio} di\nl
habil industriisto.

%%K t04-c1-10-1 p
\VUblancAlinea
10. --- La \VUgras{spozi} esas yuna e tamen li de longe\nl
esis \VUgras{fiancita.} La \VUgras{yuna muliero} \textsuperscript{(19)}, Martha, esas\nl
charmanta en sua \VUgras{blanka robo} garnisita per\nl
flori di oranjiero.

%%K t04-c1-11-1 p
\VUblancAlinea
11. --- Elua \VUgras{bopatrulo} \textsuperscript{(20)} esas tre felica havar\nl
tante jentila \VUgras{bofiliino,} e la \VUgras{bomatro} \textsuperscript{(21)} di Ioan\cc
nes ridetas vidante tante bela sua \VUgras{filiino.}

%%K t04-c1-12-1 p
\VUblancAlinea
12. --- Ye l'extremajo di la tablo, sidas la\nl
cetera \VUgras{invititi} \textsuperscript{(22)}, \VUgras{parenti, amiki, kuzuli, ku}\cc
\VUgras{zini.} \VUgras{Tablo-chefo} \textsuperscript{(23)}, kun servo-tuko sub bra\cc
kio, diligente servas li.

%%K t04-tit-4 sub
\VUpk{10.2pt}{40}{La Amiki.}
\VUsaut{3.0mm}

%%K t04-c1-13-1 p
\VUblancAlinea
13. --- Dextralatere di la tablo, yen dam\cc
\VUgras{zelo} \textsuperscript{(24)}, \VUgras{amiko} di la yuna spozino, pose la
\parplein
\end{VUpage}

% --- Feuillet 28 (folio 26) -----------------------------------------
\begin{VUpage}[28]{26}
%%K t04-c1-13-1 p suite
\VUcontinue
\VUgras{fratulo} di elca, ilqua esas elua \VUgras{honor-kom}\cc
\VUgras{pano} \textsuperscript{(25)}. Il babilas kun sua \VUgras{honor-damzelo} \textsuperscript{(26)}\nl
fratino dil spozulo e qua, per ta mariajo, dive\cc
nas la \VUgras{bofratino} di Martha. El esas charmanta\nl
yuno. El havas bela \VUgras{juveli, perlo-koliaro,} ed\nl
ora \VUgras{braceleto.}

%%K t04-c1-14-1 p
\VUblancAlinea
14. --- Proxim li, la sioro qua stacas e levas\nl
sua glaso esas \VUgras{amiko} \textsuperscript{(27)} di la familio. Il esas\nl
un de la \VUgras{testi dil spozulo.} Il \VUgras{tostas,} drinkas por\nl
la \VUgras{saneso} di la yuna spozi e deziras a li longa\nl
feliceso. Il esas vestizita per ceremonial kos\cc
tumo, per frako kun \VUgras{shui vernisizita} \textsuperscript{(28)}. Il\nl
esas l'akompananto di la \VUgras{avino} \textsuperscript{(29)} qua esas\nl
\VUgras{vidva.}

%%K t04-c1-15-1 p
\VUblancAlinea
15. --- La yunulo en \VUgras{frako} \textsuperscript{(31)}, kun dina la\cc
biobarbo nigra, esas la \VUgras{bofrato} \textsuperscript{(30)} di Ioannes.\nl
Il esas eminenta injenioro. Il esas la vicino di\nl
\VUgras{yuna damo} \textsuperscript{(33)} tre eleganta, la \VUgras{seniora fratino}\nl
di Martha e la matro di la miniona puerineto\nl
qua venas, donante la brakio a sua kuzulo,\nl
ofrar floro e \VUgras{dezirar} ad elu \VUgras{bonvespero.} Ta\nl
damo havas tre bela \VUgras{orelringi} ed eleganta \VUgras{kap}\cc
\VUgras{vesto.}

%%K t04-c1-16-1 p
\VUblancAlinea
16. --- La mikra kuzulo, tante stranja pro\nl
sua \VUgras{bluzo} \textsuperscript{(36)} e sua bufanta \VUgras{bracho} \textsuperscript{(37)} esas tre\nl
kompatinda. Il esas \VUgras{orfano} \textsuperscript{(35)}. Tre yuna il\nl
perdis sua patrulo e sua matro. Lua onklulo\nl
esas lua \VUgras{tutelero} \textsuperscript{(38)}, e restis celiba por edukar\nl
la kompatinda infanto. Il esas benigna viro.\nl
Il esas kelke kalva e tre miopa; il uzas \VUgras{(naz)}\cc
\VUgras{binoklo.}

%%K t04-tit-5 sub
\VUsaut{2.4mm}
\VUpk{10.2pt}{40}{La Desero.}
\VUsaut{3.0mm}

%%K t04-c1-17-1 p
\VUblancAlinea
17. --- Dop la festinani, sur tablo kovrita per\nl
\VUgras{(tablo)-tuko,} la \VUgras{desero} \textsuperscript{(39)} esas preparita. En
\parplein
\end{VUpage}

% --- Feuillet 29 (folio 27) -----------------------------------------
\begin{VUpage}[29]{27}
%%K t04-c1-17-1 p suite
\VUcontinue
kupego, yen frukti, \VUgras{persiki} \textsuperscript{(40)}, \VUgras{piri} \textsuperscript{(41)},\nl
\VUgras{pomi} \textsuperscript{(42)}, \VUgras{abrikoti} \textsuperscript{(43)} e \VUgras{vitberi} \textsuperscript{(44)}. Proxime,\nl
\VUgras{ananasi} \textsuperscript{(45)}, bela \VUgras{kuko} \textsuperscript{(46)}, \VUgras{liquor-boteli} \textsuperscript{(48)} e\nl
\VUgras{champanio} \textsuperscript{(49)} ed anke koloni de \VUgras{pladi} \textsuperscript{(47)} neta.

%%K t04-c1-18-1 p
\VUblancAlinea
18. --- La \VUgras{keleristo} \textsuperscript{(50)} desstopas \VUgras{botelo} \textsuperscript{(52)}\nl
de olda vino per \VUgras{korko-tirilo} \textsuperscript{(51)}. La \VUgras{korka}\nl
\VUgras{stopilo} \textsuperscript{(53)} semblas tre desfacile tirebla. Ta ke\cc
leristo havas \VUgras{servo-tuko} \textsuperscript{(54)} sub brakio.

%%K t04-c1-19-1 p
\VUblancAlinea
19. --- Pos la dineo, la siori adiros la fu\cc
meyo, e la dami iros prontigar su por la balo.

%%K t04-tit-6 sub
\VUsaut{5.0mm}
\VUcentre{11.4pt}{0}{\textit{Duesma ceno.}}
\VUsaut{1.6mm}
\VUpk{11.4pt}{40}{L'Aniversario dil avulo.}
\VUsaut{2.6mm}

%%K t04-tit-7 sub
\VUpk{10.2pt}{40}{La Vizito.}
\VUsaut{3.0mm}

%%K t04-c2-01-1 p
\VUblancAlinea
1. --- Dek yari pasis, la gepatri di Ioannes\nl
divenis oldi tre amanta sua tri nepoti, du \VUgras{ne}\cc
\VUgras{potuli} \textsuperscript{(2)} e un \VUgras{nepotino} \textsuperscript{(3)}. Pro ke cadie even\cc
tas \VUgras{l'aniversario dil avulo,} la filii venas dezirar\nl
a li bona nasko-festo.

%%K t04-c2-02-1 p
\VUblancAlinea
2. --- La mikra Petrus esas ankore tre yuna,\nl
il evas nur tri yari. Il vidas la \VUgras{kato} \textsuperscript{(1)} proxi\cc
meskar. Il volas karezar lua bela \VUgras{pilaro} silka\cc
tra e lua longa \VUgras{kaudo,} il ne reflektas, ke sub\nl
la gracioza \VUgras{pedi,} celesas maligna ungli pintoza\nl
qui forsan skrachos ilu.

%%K t04-c2-03-1 p
\VUblancAlinea
3. --- Lua fratino Gabriela, qua donas a lu la\nl
manuo, esas multe plu serioza e Marcellus kun\nl
sua granda \VUgras{mantelo} \textsuperscript{(4)} e ledra \VUgras{zono} \textsuperscript{(5)} aspektas\nl
kelke vireskinta, malgre sua kurta bracho qua\nl
lasas videsar sua \VUgras{kalzi} \textsuperscript{(6)}.
\end{VUpage}

% --- Feuillet 30 (folio 28) -----------------------------------------
\begin{VUpage}[30]{28}
%%K t04-c2-04-1 p
\VUblancAlinea
4. --- Lia patrulo metis sua dika \VUgras{surtuto} \textsuperscript{(7)}\nl
vintrala kun \VUgras{fur-kolumo} \textsuperscript{(8)} e, en sua \VUgras{po}\cc
\VUgras{sho} \textsuperscript{(9)}, il havas fulardo. Lua spozino havas\nl
sua felta chapelo garnisita per \VUgras{plumi} \textsuperscript{(10)}, sua\nl
\VUgras{jako} \textsuperscript{(11)}, sua \VUgras{boa-furo} \textsuperscript{(12)} e sua \VUgras{mufo} \textsuperscript{(13)}.

%%K t04-c2-05-1 p
\VUblancAlinea
5. --- Esas tre kolda, nam regnas la vintro.\nl
La \VUgras{nivo} \textsuperscript{(19)} faladis dum la tota jorno e la tekti\nl
di la domi esas tote blanka. Samtempe kam\nl
la \VUgras{nokto,} la koldeso divenas mem plu akra.\nl
En la cielo, la \VUgras{luno} \textsuperscript{(17)} e la \VUgras{steli} \textsuperscript{(18)} brilas e\nl
trapasas \VUgras{l'obskureso.}

%%K t04-tit-8 sub
\VUsaut{4.2mm}
\VUpk{10.2pt}{40}{La Maladeso.}
\VUsaut{3.0mm}

%%K t04-c2-06-1 p
\VUblancAlinea
6. --- La kompatinda \VUgras{blindo} \textsuperscript{(20)} qua trairas\nl
la placo avan \VUgras{l'apoteko} \textsuperscript{(16)} duktata da sua \VUgras{pu}\cc
\VUgras{delo} \textsuperscript{(21)}, esas tre desfelica. Iustinia, la \VUgras{servis}\cc
\VUgras{tino} \textsuperscript{(14)}, adportas de la koqueyo \VUgras{boledo} \textsuperscript{(15)} de\nl
tre varma \VUgras{tizano} sukrizita jeneroze.

%%K t04-c2-07-1 p
\VUblancAlinea
7. --- La \VUgras{avino} \textsuperscript{(23)} qua volas serchar sua\nl
\VUgras{orelo-binoklo} \textsuperscript{(26)} sur la tablo, por lektar la\nl
\VUgras{preskripto} \textsuperscript{(27)} quan la \VUgras{mediko} \textsuperscript{(25)} jus skribis,\nl
su levas de sua stulego, apogante su sur sua\nl
\VUgras{bastono} \textsuperscript{(24)}.

%%K t04-c2-08-1 p
\VUblancAlinea
8. --- Ofte el sufras pro \VUgras{maladetesi, vertiji,}\nl
\VUgras{katareti, dento-dolori,} lua gambi esas febla e\nl
lua okuli ne plus esas tre bona. Malgre to, vo\cc
lunte el mokas sua olda \VUgras{mediko} qua sempre\nl
volas preskriptar ad el \VUgras{pociono} \textsuperscript{(28)}. Cadie el\nl
esas tre joyoza pro ke el havas cirkum su sua\nl
yuna familio.

%%K t04-c2-09-1 p
\VUblancAlinea
9. --- Esas komforte en la chambro bone\nl
klozita. En la marmora \VUgras{kameno} \textsuperscript{(29)} flamifas
\parplein
\end{VUpage}

% --- Feuillet 31 (folio 29) -----------------------------------------
\begin{VUpage}[31]{29}
%%K t04-c2-09-1 p suite
\VUcontinue
\VUgras{belega fairo} \textsuperscript{(30)}, ed on sentas su felica en la\nl
dolca \VUgras{lumo} di la \VUgras{lampo} \textsuperscript{(31)} quan on jus acen\cc
dis.

%%K t04-tit-9 sub
\VUsaut{4.2mm}
\VUpk{10.2pt}{40}{La Avulo.}
\VUsaut{3.0mm}

%%K t04-c2-10-1 p
\VUblancAlinea
10. --- Mem la \VUgras{avulo} \textsuperscript{(33)} obliviis, ke il esis\nl
malada, ke lua \VUgras{reumatismo} fixigas lu an sua\nl
\VUgras{stulego} \textsuperscript{(34)} e ke, hiere ankore, il havis \VUgras{febro}\nl
\VUgras{e sinkopi.} Il ne plus memoras, ke dum la lasta\nl
vintro il sufris pro \VUgras{pulmonito} e ke rauka e pe\cc
nigiva \VUgras{tusado} igis lu tre sufranta.

%%K t04-c2-10-2 p
\VUblancAlinea
E tamen tre desfacile il esis risanigata. De\nl
nun il standas kelke plu bone. Ma lua gambo\nl
quan il apogas sur \VUgras{kuseno} \textsuperscript{(38)} pozita sur mi\cc
kra \VUgras{tabureto} \textsuperscript{(39)} anke dolorigas ilu.

%%K t04-c2-11-1 p
\VUblancAlinea
11. --- Kande il esas sorgoze envelopita en\nl
sua varma \VUgras{chambro-robo} \textsuperscript{(35)} kun dika \VUgras{bu}\cc
\VUgras{toni} \textsuperscript{(36)} perlomatra e kande il havas proxim\nl
su, por lektar ad ilu la jurnalo, la mikra \VUgras{orfa}\cc
\VUgras{nino} \textsuperscript{(40)} vestizita per blanka \VUgras{kofio,} per blua\nl
\VUgras{robo} \textsuperscript{(42)} e \VUgras{pelerino} \textsuperscript{(41)}, il ne plendas.

%%K t04-c2-12-1 p
\VUblancAlinea
12. --- Singlayare, kande la \VUgras{kalendario} \textsuperscript{(43)}\nl
itere indikas la \VUgras{dato} dil unesma \VUgras{decembro,} il\nl
nepaciente vartas sua \VUgras{nepoti,} nam il savas, ke\nl
li ne oblivios ta \VUgras{dato} importanta.

%%K t04-c2-13-1 p
\VUblancAlinea
13. --- Il emoce memoras la tempo tre olima\nl
kande il persone venis dezirar bona festo al\nl
glorioza \VUgras{marshalo} imperiala, di qua la por\cc
\VUgras{treto} \textsuperscript{(44)} pendas an la muro, e qua esas la\nl
\VUgras{preavo} di lua filii.

\VUsaut{6.0mm}
\VUfilet{22mm}
\end{VUpage}
